\documentclass[11pt]{article}
\usepackage[margin=1in]{geometry}
\usepackage{amsmath}

\setlength{\parindent}{0pt}

\begin{document}

\textbf{Question 6.5}
\bigskip

Suppose a new process was developed that could be used to make oil out of seawater. The equipment required is quite expensive, but it would in time lead to low prices for gasoline, electricity, and other types of energy. What effect would this have on interest rates?

\bigskip
\bigskip

The expensive equipment creates massive new investment opportunities, so firms borrow heavily to build these facilities. This shifts the demand for loanable funds to the right, pushing interest rates up. Over time, cheaper energy raises real incomes across the economy, which increases savings. This shifts the supply of loanable funds to the right, putting downward pressure on rates. In the short run, the surge in investment demand dominates, so interest rates rise. In the long run, as higher savings kick in, rates would moderate.

\bigskip\bigskip

\textbf{Question 6.8}
\bigskip

Suppose interest rates on Treasury bonds rose from 5\% to 9\% as a result of higher interest rates in Europe. What effect would this have on the price of an average company's common stock?

\bigskip
\bigskip

Stock prices would fall. Treasury bonds are a relatively riskless alternative to stocks. When Treasuries now yield 9\% instead of 5\%, investors demand a higher return from stocks to justify the extra risk. A higher required return means future cash flows are discounted at a higher rate, reducing the present value of stocks. Additionally, some investors will sell stocks and buy the now-more-attractive bonds, further pushing stock prices down.

\bigskip\bigskip

\textbf{Problem 6.3}
\bigskip

The real risk-free rate is 2.25\%. Inflation is expected to be 2.5\% this year and 4.25\% during the next 2 years. Assume that the maturity risk premium is zero. What is the yield on 2-year Treasury securities? What is the yield on 3-year Treasury securities?

\bigskip
\bigskip

2-year Treasury:
\[
\text{IP}_2 = \frac{2.5\% + 4.25\%}{2} = 3.375\%
\]
\[
r_2 = 2.25\% + 3.375\% = \boxed{5.625\%}
\]

3-year Treasury:
\[
\text{IP}_3 = \frac{2.5\% + 4.25\% + 4.25\%}{3} \approx 3.667\%
\]
\[
r_3 = 2.25\% + 3.667\% = \boxed{5.917\%}
\]

\bigskip\bigskip

\textbf{Problem 6.5}
\bigskip

The real risk-free rate is 2.5\% and inflation is expected to be 2.75\% for the next 2 years. A 2-year Treasury security yields 5.55\%. What is the maturity risk premium for the 2-year security?

\bigskip
\bigskip

\[
\text{IP}_2 = \frac{2.75\% + 2.75\%}{2} = 2.75\%
\]
\[
5.55\% = 2.5\% + 2.75\% + \text{MRP}
\]
\[
\text{MRP} = 5.55\% - 2.5\% - 2.75\% = \boxed{0.30\%}
\]

\bigskip\bigskip

\textbf{Additional Question}
\bigskip

A bond trader observes: the Treasury yield curve is downward sloping; a positive maturity risk premium applies to both Treasury and corporate bonds; there is no liquidity premium for Treasuries but a positive liquidity premium is built into corporate bond yields. Compare a 5-year corporate bond and a 10-year Treasury bond. Which must have a higher yield?

\bigskip
\bigskip

The 5-year corporate bond must have the higher yield. Since the Treasury yield curve is downward sloping, the 5-year Treasury yields more than the 10-year Treasury (declining inflation expectations more than offset the positive MRP). The 5-year corporate bond yields even more than the 5-year Treasury because it adds a default risk premium and a liquidity premium on top. So: 5-year corporate $>$ 5-year Treasury $>$ 10-year Treasury. The 5-year corporate bond must have the higher yield.

\end{document}
